% Créditos -------------------------------------------------------------------

% Algunas opciones de configuración fueron tomadas y adaptadas del estilo
% del libro IMS en pdf
% https://github.com/OpenIntroStat/ims/blob/main/latex/ims-style.tex

% PAQUETES -------------------------------------------------------------------

\usepackage{fontspec}
\usepackage{colortbl, xcolor}
\usepackage[explicit]{titlesec}     % para formato de titulos
\usepackage{fancyhdr}
\usepackage{lipsum}
\usepackage{tikz}
\usepackage[framemethod=tikz]{mdframed}  % para ambientes

% VARIAS CONFIGURACIONESS  ----------------------------------------------------

% Márgenes 
\usepackage[top=3cm,bottom=3cm,left=3cm,right=1.5cm]{geometry}
 
% Colores
\definecolor{oiB}{HTML}{569BBD}            % COL["blue","full"]
\definecolor{oiLB}{HTML}{e3eef4}           % lighter version of oiB

\definecolor{oiY}{HTML}{f4dc00}            % COL["yellow","full"]
\definecolor{oiLY}{HTML}{fffacd}           % lighter version of oiY

\definecolor{oiR}{HTML}{E97583}            % COL["red","full"]
\definecolor{oiLR}{HTML}{F3CED4}           % lighter version of oiR

\definecolor{oiGray}{HTML}{808080}         % COL["gray","full"]
\definecolor{oiLGray}{HTML}{f8f8f8}        % lighter version of oiGray

\definecolor{oiGB}{rgb}{0.5,0.5,.5}        % from OS4 - for footnotes

% TOC
% que no tenga el nro de parte, porque eso se lo pongo yo en el título
% Deshabilitar la numeración en la tabla de contenidos
% \renewcommand{\thepart}{\relax}
% ESO NO, LO SACA DE TODOS LADOS

% % Redefinir el comando \part para evitar la carátula 
% esto hace que el texto de la parte este a continuacion del titulo
\titleclass{\part}{top}
% esto modifica el formato del titulo
\titleformat{\part}[display]
{\color{oiR}\titlerule[5pt]\vspace{3pt}\color{oiLR}\titlerule[2pt]\color{oiB}\normalfont\Huge\bfseries\scshape}
{}{0em}{#1 \\ \noindent \vspace{3pt}\color{oiR}\titlerule[5pt]}
\titlespacing*{\part}{0pt}{*0}{15em}

% Encabezado y pie de página --------------------------------------------------

% Paquete para personalizar encabezados y pies de página
% Configuración de fancyhdr
\pagestyle{fancy}

% Limpiar estilos predeterminados
\fancyhf{}

% Configuración del contenido
\fancyhead[LO, RE]{Unidad \thepart}
\fancyhead[RO, LE]{\nouppercase{\leftmark}}
\fancyfoot[LO, RE]{Programación 1}
\fancyfoot[RO, LE]{\thepage}

\renewcommand{\footrulewidth}{0.4pt}
\renewcommand{\sectionmark}[1]{\markright{#1}}

% para que la primera pagina no tenga el header pero si el footer como el resto
\fancypagestyle{primerapagina}{
  \fancyfoot[LO, RE]{Programación 1}
  \fancyfoot[RO, LE]{\thepage}
	\fancyhead{}
	\renewcommand{\headrulewidth}{0pt}
}

% Formato de títulos ----------------------------------------------------------

\titleformat{\chapter}[display]
{\color{oiB}\normalfont\Huge\bfseries}
{\color{oiB}Capítulo \thechapter}{1em}{#1}

\titleformat{\section}
{\color{oiB}\normalfont\Large\bfseries}
{\color{oiB}\thesection}{1em}{#1}

\titleformat{\subsection}
{\color{oiB}\normalfont\large\bfseries}
{\color{oiB}\thesubsection}{1em}{#1}

\titleformat{\subsubsection}
{\color{oiB}\normalfont\normalsize\bfseries}
{\color{oiB}\thesubsubsection}{1em}{#1}


% Environments -----------------------------------------------------------

% Helper environments
% currently none of them increment footnote counter,
% this seems to be the right approach?

% mdframedwithfootChapterintro: for presentacion-capitulo
\newenvironment{mdframedwithfootChapterintro}
{   
    \savenotes
    \begin{mdframed}[%
    topline=true, bottomline=true, linecolor=oiB, linewidth=1.4pt,
    rightline=false, leftline=false,
    backgroundcolor=oiLB]
    %\stepcounter{footnote} % don't increment footnote counter
    \renewcommand{\thempfootnote}{\arabic{footnote}}
    }
{
    \end{mdframed}
    \spewnotes
}

% mdframedwithfootGPWE: for ejercicio
\newenvironment{mdframedwithfootGPWE}
{   
    \savenotes
    \begin{mdframed}[%
    topline=true, bottomline=true, linecolor=oiB, linewidth=0.5pt,
    rightline=false, leftline=false,
    backgroundcolor=oiLGray]
    %\stepcounter{footnote}
    \renewcommand{\thempfootnote}{\arabic{footnote}}
    }
{
    \end{mdframed}
    \spewnotes
}

% mdframedwithfootImportant: for ejemplo
\newenvironment{mdframedwithfootImportant}
{   
    \savenotes
    \begin{mdframed}[%
    topline=true, bottomline=true, linecolor=oiR, linewidth=0.5pt,
    rightline=false, leftline=false,
    backgroundcolor=oiLGray]
    %\stepcounter{footnote}
    \renewcommand{\thempfootnote}{\arabic{footnote}}
    }
{
    \end{mdframed}
    \spewnotes
}

% mdframedwithfootTip: for concepto, info, comentario, en-contruccion
\newenvironment{mdframedwithfootTipDataPro}
{   
    \savenotes
    \begin{mdframed}[%
    topline=true, bottomline=true, linecolor=oiGray, linewidth=0.5pt,
    rightline=false, leftline=false,
    backgroundcolor=oiLGray]
    %\stepcounter{footnote}
    \renewcommand{\thempfootnote}{\arabic{footnote}}
    }
{
    \end{mdframed}
    \spewnotes
}

% los ambientes uses regular mdframedwithfoot variants, might have footnotes

% Intro capitulo
\newenvironment{presentacion-capitulo}{
\vspace{4mm}
\begin{mdframedwithfootChapterintro}
\begin{minipage}[t]{0.10\textwidth}
{$\:$ \\ \setkeys{Gin}{width=2.5em,keepaspectratio}\includegraphics{imagenes/iconos/006-clean-code.png}}
\end{minipage}
\hfill
\begin{minipage}[t]{0.90\textwidth}
\setlength{\parskip}{1em}
\noindent\textbf{\color{oiB}\small{\MakeUppercase{Resumen}}} $\:$ \\ 
\large
}{\end{minipage}
\end{mdframedwithfootChapterintro}
\vspace{4mm}
}

% Ejercicio
\newenvironment{ejercicio}{
\vspace{4mm}
\begin{mdframedwithfootGPWE}
\begin{minipage}[t]{0.10\textwidth}
{$\:$ \\ \setkeys{Gin}{width=2.5em,keepaspectratio}\includegraphics{imagenes/iconos/008-edit-code.png}}
\end{minipage}
\hfill
\begin{minipage}[t]{0.90\textwidth}
%\vspace{-2mm}
\setlength{\parskip}{1em}
\noindent\textbf{\color{oiB}\small{\MakeUppercase{Para resolver}}} $\:$ \\ 
}{\end{minipage}
\end{mdframedwithfootGPWE}
\vspace{4mm}
}

% ejemplo
\newenvironment{ejemplo}{
    \let\oldrule\rule
    \renewcommand{\rule}[2]{\vspace{-2mm}\oldrule{##1}{##2}\vspace{-2mm}}
\vspace{4mm}
\begin{mdframedwithfootImportant}
\begin{minipage}[t]{0.10\textwidth}
{$\:$ \\ \setkeys{Gin}{width=2.5em,keepaspectratio}\includegraphics{imagenes/iconos/code-review.png}}
\end{minipage}
\hfill
\begin{minipage}[t]{0.90\textwidth}
%\vspace{-2mm}
\setlength{\parskip}{1em}
\noindent\textbf{\color{oiB}\small{\MakeUppercase{Ejemplo}}} $\:$ \\
\small}{\end{minipage}
\end{mdframedwithfootImportant}
\vspace{4mm}
}

% concepto
\newenvironment{concepto}{
    \let\oldtextbf\textbf
    \renewcommand{\textbf}[1]{{\textcolor{oiR}{\oldtextbf{##1}}}}
\vspace{4mm}
\begin{mdframedwithfootTipDataPro}
\begin{minipage}[t]{0.10\textwidth}
{$\:$ \\ \setkeys{Gin}{width=2.5em,keepaspectratio}\includegraphics{imagenes/iconos/001-anchor.png}}
\end{minipage}
\hfill
\begin{minipage}[t]{0.90\textwidth}
%\vspace{-2mm}
\setlength{\parskip}{1em}
\noindent\textbf{\color{oiB}\small{\MakeUppercase{Concepto clave}}} $\:$ \\
\small}{\end{minipage}
\end{mdframedwithfootTipDataPro}
\vspace{4mm}
}


% info
\newenvironment{info}{
    \let\oldtextbf\textbf
    \renewcommand{\textbf}[1]{{\textcolor{oiR}{\oldtextbf{##1}}}}
\vspace{4mm}
\begin{mdframedwithfootTipDataPro}
\begin{minipage}[t]{0.10\textwidth}
{$\:$ \\ \setkeys{Gin}{width=2.5em,keepaspectratio}\includegraphics{imagenes/iconos/010-review.png}}
\end{minipage}
\hfill
\begin{minipage}[t]{0.90\textwidth}
%\vspace{-2mm}
\setlength{\parskip}{1em}
\noindent\textbf{\color{oiB}\small{\MakeUppercase{Info imporante}}} $\:$ \\
\small}{\end{minipage}
\end{mdframedwithfootTipDataPro}
\vspace{4mm}
}


% comentario DESPRECIADO PORQUE SE ROMPE SI ES MAS LARGO DE UNA PAGINA
% \newenvironment{comentario}{
%     \let\oldtextbf\textbf
%     \renewcommand{\textbf}[1]{{\textcolor{oiR}{\oldtextbf{##1}}}}
% \vspace{4mm}
% \begin{mdframedwithfootTipDataPro}
% \begin{minipage}[t]{0.10\textwidth}
% {$\:$ \\ \setkeys{Gin}{width=2.5em,keepaspectratio}\includegraphics{imagenes/iconos/007-browser.png}}
% \end{minipage}
% \hfill
% \begin{minipage}[t]{0.90\textwidth}
% %\vspace{-2mm}
% \setlength{\parskip}{1em}
% \noindent\textbf{\color{oiB}\small{\MakeUppercase{Comentario adicional}}} $\:$ \\
% \small}{\end{minipage}
% \end{mdframedwithfootTipDataPro}
% \vspace{4mm}
% }

% en-construccion
\newenvironment{en-construccion}{
    \let\oldtextbf\textbf
    \renewcommand{\textbf}[1]{{\textcolor{oiR}{\oldtextbf{##1}}}}
\vspace{4mm}
\begin{mdframedwithfootTipDataPro}
\begin{minipage}[t]{0.10\textwidth}
{$\:$ \\ \setkeys{Gin}{width=2.5em,keepaspectratio}\includegraphics{imagenes/iconos/003-website.png}}
\end{minipage}
\hfill
\begin{minipage}[t]{0.90\textwidth}
%\vspace{-2mm}
\setlength{\parskip}{1em}
\noindent\textbf{\color{oiB}\small{\MakeUppercase{En desarrollo}}} $\:$ \\
\small}{\end{minipage}
\end{mdframedwithfootTipDataPro}
\vspace{4mm}
}

%---------------
% NUEVOS AMBIENTES QUE SE EXTIENDE A TRAVES DE VARIAS PAGINAS

\usepackage[most]{tcolorbox}


\newtcolorbox{comentario}{
  breakable,
  enhanced,
  fonttitle=\sffamily\bfseries,
  coltitle=oiB,
  title=\textbf{\color{oiB}\small{\MakeUppercase{Comentario adicional}}} $\:$,
  title code={
  \path[fill=oiLGray] (title.south west) rectangle (title.north east);
  \path[draw=oiB,solid,line width=0.75mm]
  ([xshift=5mm]title.south west) -- ([xshift=-5mm]title.south east);
  },
  boxrule=0pt,frame hidden,
  borderline west={3pt}{0pt}{oiB},
  colback=oiLGray,
  colbacktitle=oiLGray,
}

